% ============================================================================
\chapter{Taller práctico: 8 ejercicios guiados}\label{cap:taller}
% ============================================================================

En este taller hará, una vez y de principio a fin, todo lo que hará después en su trabajo real. Cada ejercicio indica
\textbf{qué escribir y en qué celda}, y al final trae un cuadro verde con \textbf{los números exactos que debe ver}.
Si coinciden, lo hizo bien.

\begin{atencion}
\textbf{Haga los ejercicios sobre una copia de práctica} (sección~\ref{sec:practica}): \texttt{PRACTICA\_Rentabilidad.xlsx}.
Hágalos \textbf{en orden}: cada ejercicio parte de cómo quedó el anterior. Si se equivoca, use \tecla{Ctrl}+\tecla{Z}.
Si se pierde del todo, vuelva a copiar el archivo original y empiece de nuevo.
\end{atencion}

\begin{center}
\begin{tabularx}{\linewidth}{@{}clX@{}}
\toprule
\textbf{Nº} & \textbf{Ejercicio} & \textbf{Qué aprende} \\ \midrule
1 & Comprobar el estado del libro & Leer los controles de \hoja{INICIO} \\
2 & Registrar un contenedor que llega & El primer momento de cada contenedor \\
3 & Completar los datos reales y cerrarlo & El segundo momento; cómo nace una alerta \\
4 & Sube el precio del diésel & Agregar una vigencia de precios sin alterar el pasado \\
5 & Cierre del mes de octubre & Gastos reales, liquidación de los gestores y conciliación \\
6 & Informe del período & Usar el \hoja{DASHBOARD} \\
7 & ¿Y si subimos la tarifa de Oriente? & Usar el \hoja{SIMULADOR} \\
8 & Dejar el libro limpio para empezar & Borrar los ejemplos de forma segura \\
\bottomrule
\end{tabularx}
\end{center}

% ============================================================================
\section{Ejercicio 1. Comprobar el estado del libro}

\begin{ficha}{Ejercicio 1 · Duración: 3 minutos}
\textbf{Objetivo:} aprender a leer el tablero de control antes de trabajar.\\
\textbf{Hoja:} \hoja{INICIO}
\end{ficha}

\begin{pasos}
\paso Abra \texttt{PRACTICA\_Rentabilidad.xlsx}. Si aparece la barra amarilla de \textbf{Vista protegida}, pulse \boton{Habilitar edición}.
\paso Haga clic en la pestaña \hoja{INICIO}.
\paso Baje hasta el bloque \textbf{Estado del libro (controles automáticos)}, alrededor de la fila 35.
\paso Lea cada línea y compárela con el cuadro verde.
\end{pasos}

\captura{inicio_controles}{Ejercicio 1: el estado inicial del libro}

\begin{resultado}
Contenedores registrados \textbf{4}, con alertas \textbf{\vAlertasIni}, abiertos \textbf{\vAbiertosIni};
cuadres \textbf{OK} y \textbf{OK}; conciliación \textbf{CUADRA}; utilidad acumulada \textbf{\vUAItotal}.
\end{resultado}

\begin{nota}
Haga clic en algún nombre de hoja del índice (filas 10 a 22): Excel lo lleva a esa hoja. Para volver, haga clic en la pestaña \hoja{INICIO}.
\end{nota}

% ============================================================================
\section{Ejercicio 2. Registrar un contenedor que llega}\label{sec:ej2}

\begin{ficha}{Ejercicio 2 · Duración: 5 minutos}
\textbf{Situación:} el 27 de octubre de 2026 llega el contenedor \textbf{CONT-PRUEBA-05}. Pasa por ENCI, es un 40' HC con 1\,000 bultos.
Según el manifiesto, lleva \textbf{5\,000 kg} para Occidente, \textbf{2\,500 kg} para Centro y \textbf{2\,500 kg} para Oriente.\\
\textbf{Hoja:} \hoja{ENTRADA} · \textbf{Fila:} 10, la primera vacía
\end{ficha}

\begin{pasos}
\paso Vaya a la hoja \hoja{ENTRADA}. Las filas 6 a 9 tienen los 4 contenedores de ejemplo. La primera fila libre es la \textbf{10}.
\paso Haga clic en \celda{B10} y escriba \escriba{CONT-PRUEBA-05}. Pulse \tecla{Tab}: pasa a \celda{C10}. En \celda{A10} aparece el número 5.
\paso En \celda{C10} escriba la fecha \escriba{27/10/2026} y pulse \tecla{Tab}.
\paso En \celda{D10} pulse la flecha de la lista y elija \textbf{ENCI}. Pulse \tecla{Tab}.
\paso En \celda{E10} elija \textbf{40' HC} de la lista. Pulse \tecla{Tab}.
\paso En \celda{F10} escriba \escriba{1000} y pulse \tecla{Tab}.
\paso En \celda{G10} escriba \escriba{5000}; en \celda{H10}, \escriba{2500}; en \celda{I10}, \escriba{2500}. Pulse \tecla{Tab} entre una y otra.
\paso Deje vacías las columnas J a X: todavía no tiene los datos reales.
\paso Vaya a \celda{Z10} (pulse \tecla{F5}, escriba \escriba{Z10} y \tecla{Enter}) y elija \textbf{Abierto} en la lista.
\paso Guarde con \tecla{Ctrl}+\tecla{G}.
\end{pasos}

\captura{ej2_entrada}{Ejercicio 2: el contenedor nuevo en la fila 10 (columnas A a I)}
\captura{ej2_estado}{Ejercicio 2: estado y resultado del contenedor (columnas Y a AC)}

\begin{resultado}
\begin{itemize}
  \item \recuadro{1} \celda{Z10} dice \textbf{Abierto}.
  \item \recuadro{2} \celda{AB10} (alertas) está \textbf{vacía} y \celda{AC10} (utilidad) muestra \textbf{\vEjDosUAI}.
\end{itemize}
\end{resultado}

\begin{nota}
\textbf{¿Por qué cambió la utilidad de otro contenedor?} Antes del ejercicio, el contenedor 3 (fila 8) tenía una utilidad de \vUAIcTres.
Ahora tiene \vEjDosUAItres. No es un error: los \textbf{gastos fijos de octubre} se reparten entre todos los contenedores de octubre según
sus kilogramos. Al llegar uno más, a cada uno le toca una parte menor. Por eso \textbf{el resultado de un mes es definitivo solo cuando el
mes se cierra}.
\end{nota}

% ============================================================================
\section{Ejercicio 3. Completar los datos reales y cerrar el contenedor}\label{sec:ej3}

\begin{ficha}{Ejercicio 3 · Duración: 5 minutos}
\textbf{Situación:} terminó la distribución del CONT-PRUEBA-05 el 6 de noviembre. Los vales dicen que el viaje al puerto consumió
\textbf{64 litros}, que el desagrupe en ENCI llevó \textbf{11 jornales} y que el traslado a Oriente consumió \textbf{520 litros}.
Los demás datos fueron los normales.\\
\textbf{Hoja:} \hoja{ENTRADA} · \textbf{Fila:} 10
\end{ficha}

\begin{pasos}
\paso En \celda{K10} (Litros reales al puerto) escriba \escriba{64}.
\paso En \celda{M10} (Jornales del desagrupe) escriba \escriba{11}.
\paso En \celda{S10} (Litros de traslado a Oriente) escriba \escriba{520}.
\paso Las demás columnas opcionales quedan vacías: el libro usará la norma.
\paso En \celda{Y10} escriba la fecha de fin \escriba{06/11/2026}.
\paso En \celda{Z10} cambie el estado a \textbf{Cerrado}.
\paso En \celda{AA10} anote los documentos: \escribal{Vales 1520-1528; nómina de desagrupe 45}.
\end{pasos}

\captura{ej3_reales}{Ejercicio 3: datos reales en las columnas opcionales}
\captura{ej3_estado}{Ejercicio 3: contenedor cerrado, con una alerta}

\begin{resultado}
\begin{itemize}
  \item \recuadro{2} \celda{AB10} muestra la alerta: «\textbf{\vEjTresAlerta}».
  \item \celda{AC10}: la utilidad baja a \textbf{\vEjTresUAI}.
  \item En \hoja{INICIO}, «Contenedores con alertas» pasa a \textbf{\vEjTresAlertas}.
\end{itemize}
\end{resultado}

\captura{ej3_inicio}{Ejercicio 3: el control de INICIO cuenta la nueva alerta}

\begin{nota}
\textbf{¿Qué significa la alerta?} La norma para el traslado a Oriente es de 450 litros por viaje. Se gastaron 520, un 16\,\% más, y la
tolerancia es del 10\,\%. El libro no impide registrarlo: \textbf{avisa} para que usted averigüe la causa (un desvío, una avería, una ruta
distinta, un vale mal anotado\dots).
\begin{itemize}
  \item Si el dato es correcto y está justificado, explíquelo en Observaciones.
  \item Si era un error de anotación, corríjalo. La alerta desaparece sola.
  \item Si el exceso se repite con todos los viajes, la norma está desactualizada: corríjala en \hoja{PARAMETROS} (capítulo~\ref{cap:cambios}).
\end{itemize}
\end{nota}

% ============================================================================
\section{Ejercicio 4. Sube el precio del diésel}\label{sec:ej4}

\begin{ficha}{Ejercicio 4 · Duración: 8 minutos}
\textbf{Situación:} desde el \textbf{1 de noviembre de 2026} el diésel pasa de 2,00 a \textbf{2,30 USD/L}. El 10 de noviembre llega el
contenedor \textbf{CONT-PRUEBA-06} por A.V: un 40' con 820 bultos y 4\,000 / 2\,000 / 2\,000 kg.\\
\textbf{Hojas:} \hoja{PRECIOS} y \hoja{ENTRADA}
\end{ficha}

\begin{atencion}
\textbf{No cambie el 2,00 de la fila 6 por 2,30.} Si lo hiciera, el libro recalcularía todos los contenedores anteriores con el precio nuevo
y el historial quedaría falseado. Lo correcto es \textbf{agregar una fila} con la nueva fecha.
\end{atencion}

\begin{pasos}
\paso Vaya a \hoja{PRECIOS}. La fila 6 tiene los precios vigentes desde el 01/01/2026.
\paso Seleccione \celda{B6:N6}: haga clic en \celda{B6}, mantenga \tecla{Mayús} y haga clic en \celda{N6}. Cópielo con \tecla{Ctrl}+\tecla{C}.
\paso Haga clic en \celda{B7} y pegue con \tecla{Ctrl}+\tecla{V}. La fila 7 queda igual que la 6.
\paso En \celda{A7} escriba \escriba{01/11/2026}: la fecha desde la que rige el cambio.
\paso En \celda{E7} (Precio del diésel) escriba \escriba{2.30}.
\paso En \celda{O7} anote el soporte: \escribal{Aumento del diésel (factura CUPET 0457)}.
\paso Pulse \tecla{Esc} para quitar la marca de copiado.
\paso Vaya a \hoja{ENTRADA} y registre el nuevo contenedor en la fila \textbf{11}, como en el ejercicio 2:
      \celda{B11} \escriba{CONT-PRUEBA-06}, \celda{C11} \escriba{10/11/2026}, \celda{D11} \textbf{A.V}, \celda{E11} \textbf{40'},
      \celda{F11} \escriba{820}, \celda{G11} \escriba{4000}, \celda{H11} \escriba{2000}, \celda{I11} \escriba{2000}, \celda{Z11} \textbf{Abierto}.
\end{pasos}

\captura{ej4_precios}{Ejercicio 4: la nueva vigencia en la fila 7 (se ocultaron las columnas G a N)}

Para comprobar que cada contenedor tomó su precio, mire la hoja \hoja{CALCULO}. La imagen muestra solo algunas columnas:

\captura{ej4_calculo}{Ejercicio 4: cada contenedor usa los precios de su fecha}

\begin{resultado}
\begin{itemize}
  \item Columna «Precios vigentes desde»: el CONT-PRUEBA-06 (fila 11) usa la vigencia del \textbf{01/11/2026}. Todos los demás siguen con la del 01/01/2026.
  \item Combustible al puerto del CONT-PRUEBA-06: \textbf{\vEjCuatroCombSeis} (60 L $\times$ 2,30). El CONT-PRUEBA-05 sigue con \textbf{\vEjCuatroCombCinco}
        (64 L $\times$ 2,00).
  \item La utilidad del CONT-PRUEBA-05 \textbf{no cambió}: sigue en \textbf{\vEjCuatroUAIcinco}. El historial está intacto.
  \item El CONT-PRUEBA-06 muestra una pérdida de \textcolor{red}{\textbf{\vEjCuatroUAIseis}}.
\end{itemize}
\end{resultado}

\begin{nota}
\textbf{¿Por qué pierde el CONT-PRUEBA-06?} Es, por ahora, el \textbf{único} contenedor de noviembre, así que carga solo con todos los gastos
fijos del mes. Cuando se registren los demás contenedores de noviembre, el gasto se repartirá y su resultado mejorará.
Un mes con pocos contenedores es caro: lo verá en el indicador «punto de equilibrio» (capítulo~\ref{cap:resultados}).
\end{nota}

% ============================================================================
\section{Ejercicio 5. Cierre del mes de octubre}\label{sec:ej5}

\begin{ficha}{Ejercicio 5 · Duración: 15 minutos}
\textbf{Situación:} es 1 de noviembre y toca cerrar octubre. La factura de electricidad de octubre fue de \textbf{55 USD}, no de 40.
Los gestores informan lo que cobraron en octubre. Al cobro habitual se suma el del CONT-PRUEBA-05:
\textbf{4\,000 USD} en La Habana, \textbf{2\,000} en Villa Clara y \textbf{2\,000} en Santiago de Cuba.\\
\textbf{Hojas:} \hoja{COSTOS\_FIJOS}, \hoja{COMERCIAL}, \hoja{RESUMEN\_MES}
\end{ficha}

\subsection*{Parte A: registrar un gasto fijo real distinto del habitual}
\begin{pasos}
\paso Vaya a \hoja{COSTOS\_FIJOS}. Busque la fila \textbf{10, «Electricidad y agua»}.
\paso Las columnas de los meses empiezan en G (sep-2026). Octubre es la columna \textbf{H}.
\paso Haga clic en \celda{H10}. Verá que contiene \texttt{=\$F10}: copia el valor base. Escriba encima \escriba{55} y pulse \tecla{Enter}.
\end{pasos}

\captura{ej5_costos}{Ejercicio 5A: gasto real de electricidad de octubre}

\begin{resultado}
En \hoja{RESUMEN\_MES}, «Costos fijos del mes» de octubre pasa de \vEjCincoFijosOctAntes{} a \textbf{\vEjCincoFijosOct}. Solo cambia octubre: los demás meses siguen con 40.
\end{resultado}

\subsection*{Parte B: liquidar y conciliar a los gestores}
\begin{pasos}
\paso Vaya a \hoja{COMERCIAL}. Compruebe que en \celda{C4} (Mes a liquidar) dice \textbf{oct-2026}. Si no, escriba \escriba{01/10/2026}.
\paso Baje a la \textbf{conciliación} (filas 35 a 40). Dice \textcolor{rojo}{\textbf{REVISAR}}, con una diferencia de
      \textcolor{red}{\vEjCincoDif}: los gestores declararon \vEjCincoDeclarada, pero los contenedores registrados suman \vEjCincoReal.
      Falta declarar el cobro del CONT-PRUEBA-05.
\end{pasos}

\captura{ej5_comercial_revisar}{Ejercicio 5B: la conciliación detecta una diferencia}

\begin{pasos}
\setcounter{paso}{2}
\paso Columna G, base variable del mes. Sume el cobro nuevo al que ya tenía cada gestor:
  \begin{itemize}
    \item La Habana, fila 19: \celda{G19} tenía \vEjCincoHabAntes. Escriba \escriba{\vEjCincoHab}.
    \item Villa Clara, fila 24: \celda{G24} tenía \vEjCincoVcAntes. Escriba \escriba{\vEjCincoVc}.
    \item Santiago de Cuba, fila 31: \celda{G31} tenía \vEjCincoScAntes. Escriba \escriba{\vEjCincoSc}.
  \end{itemize}
\end{pasos}

\begin{consejo}
Para sumar sin calculadora, escriba la operación en la celda: \escriba{=\vEjCincoHabAntes+4000}. Si su Excel usa coma decimal, escriba la
cifra con coma.
\end{consejo}

\captura{ej5_comercial_ok}{Ejercicio 5B: bases corregidas y conciliación cuadrada}

\begin{resultado}
\recuadro{4} La conciliación dice \textcolor{verde}{\textbf{CUADRA}}. El total a pagar a los 21 gestores pasa de \vEjCincoPagoTotalAntes{} a
\textbf{\vEjCincoPagoTotal}. Esa es la liquidación de octubre.
\end{resultado}

\subsection*{Parte C: revisar el resultado y guardar el respaldo}
\begin{pasos}
\setcounter{paso}{3}
\paso Vaya a \hoja{RESUMEN\_MES} y mire la fila de octubre.
\paso Vaya a \hoja{INICIO} y revise el estado del libro.
\paso Guarde el archivo y haga la copia de respaldo del mes (sección~\ref{sec:respaldo}).
\end{pasos}

\captura{ej5_resumen}{Ejercicio 5C: resultado de octubre después del cierre}

\begin{resultado}
\recuadro{1} Costos fijos de octubre: \textbf{\vEjCincoFijosOct}. \recuadro{2} Utilidad antes de impuesto de octubre: \textbf{\vEjCincoUAIoct}.
Noviembre, con un solo contenedor, va en \textcolor{red}{\vEjCincoUAInov}.
\end{resultado}

% ============================================================================
\section{Ejercicio 6. Informe del período}

\begin{ficha}{Ejercicio 6 · Duración: 3 minutos}
\textbf{Situación:} la dirección pide el resultado de octubre y noviembre.\\
\textbf{Hoja:} \hoja{DASHBOARD}
\end{ficha}

\begin{pasos}
\paso Vaya a \hoja{DASHBOARD}.
\paso Haga clic en \celda{C4} (Período desde), pulse la flecha y elija \textbf{oct-2026}. También puede escribir \escriba{01/10/2026}.
\paso En \celda{C5} (Período hasta) elija \textbf{nov-2026}.
\end{pasos}

\captura{ej6_dashboard}{Ejercicio 6: indicadores de octubre y noviembre}

\begin{resultado}
Contenedores: \textbf{\vEjSeisCont}. Utilidad antes de impuesto: \textbf{\vEjSeisUAI}. Margen neto: \textbf{\vEjSeisMargen}.
Los gráficos de más abajo se actualizan solos.
\end{resultado}

\begin{consejo}
Para presentar el informe: \menu{Archivo $\rightarrow$ Imprimir} con la hoja \hoja{DASHBOARD} activa, o \menu{Archivo $\rightarrow$
Exportar $\rightarrow$ PDF}. La hoja ya está preparada para salir en horizontal y ajustada al ancho de la página.
\end{consejo}

% ============================================================================
\section{Ejercicio 7. ¿Y si subimos la tarifa de Oriente?}

\begin{ficha}{Ejercicio 7 · Duración: 5 minutos}
\textbf{Situación:} Oriente casi no deja ganancia. La dirección pregunta qué pasaría si su tarifa subiera de 0,80 a \textbf{0,95 USD/kg}.\\
\textbf{Hoja:} \hoja{SIMULADOR}
\end{ficha}

\begin{pasos}
\paso Vaya a \hoja{SIMULADOR}. Mire la fila 11, «Tarifa Oriente».
\paso En \celda{D11} (Valor a simular) escriba \escriba{0.95}.
\paso Mire el resultado más abajo, en las filas 38 a 46.
\paso Cuando termine, \textbf{borre \celda{D11}} con \tecla{Supr} para volver al valor real.
\end{pasos}

\captura{ej7_simulador}{Ejercicio 7: tarifa de Oriente simulada}
\captura{ej7_resultado}{Ejercicio 7: nuevo resultado}

\begin{resultado}
\begin{itemize}
  \item Utilidad de Oriente por contenedor: de \vSimUAIori{} a \textbf{\vEjSieteUAIori}. Margen de Oriente: de \vSimMgOri{} a \textbf{\vEjSieteMgOri}.
  \item Utilidad del mes: de \vSimUAImes{} a \textbf{\vEjSieteUAImes}.
\end{itemize}
\end{resultado}

\begin{nota}
El simulador \textbf{nunca} modifica los datos reales. Puede probar tantos escenarios como quiera. Recuerde borrar los valores simulados al terminar
para no confundirse la próxima vez.
\end{nota}

% ============================================================================
\section{Ejercicio 8. Dejar el libro limpio para empezar}\label{sec:ej8}

\begin{ficha}{Ejercicio 8 · Duración: 5 minutos}
\textbf{Objetivo:} aprender a borrar datos sin dañar el libro. Es lo que hará en el archivo real antes de cargar sus datos.\\
\textbf{Hojas:} \hoja{ENTRADA}, \hoja{COMERCIAL}
\end{ficha}

\begin{pasos}
\paso Para el ejercicio, cierre la copia de práctica sin guardar y vuelva a hacer una copia del original: así parte de los 4 contenedores de ejemplo.
\paso En \hoja{ENTRADA}, haga clic en \celda{B6} y arrastre hasta \celda{AA9}, de modo que queden seleccionadas las 4 filas de ejemplo.
\paso Pulse \tecla{Supr}. \textbf{No} use «Eliminar filas».
\paso En \hoja{COMERCIAL}, seleccione \celda{G12:G32} (base variable de los gestores) y pulse \tecla{Supr}.
\paso Vaya a \hoja{INICIO} y mire el estado del libro.
\end{pasos}

\captura{ej8_vacio}{Ejercicio 8: ENTRADA sin los ejemplos}
\captura{ej8_inicio}{Ejercicio 8: el estado del libro vacío}

\begin{resultado}
Contenedores registrados \textbf{0}, alertas \textbf{0}, cuadres \textbf{OK}, conciliación \textbf{\vEjOchoConc} y utilidad acumulada
\textcolor{red}{\textbf{\vEjOchoUAI}}.
\end{resultado}

\begin{nota}
\textbf{¿Por qué la utilidad es negativa si no hay contenedores?} El período empieza en septiembre de 2026 (\texttt{Fecha\_Inicio}), y ese mes ya
tiene gastos fijos (\vEjOchoFijos) aunque no se haya registrado ningún contenedor. Es lo correcto: los gastos fijos se pagan haya o no carga.
Al cargar sus datos reales pondrá su propia fecha de inicio (capítulo~\ref{cap:preparar}).
\end{nota}

\begin{consejo}
\textbf{¡Felicidades!} Si llegó hasta aquí con todos los cuadros verdes coincidiendo, ya sabe hacer todo el trabajo del libro.
El siguiente paso es prepararlo con los datos reales de la empresa.
\end{consejo}
